% Generated by scripts/blueprint_gen.py; do not edit.
\chapter{Cohesive modules over a dga}\label{chap:A}
\textbf{Setting for all of §2.} k a commutative ring. A a ℤ-graded k-algebra with A\textsuperscript{q} = 0 for q < 0 and q > N, graded-commutative (ab = (−1)\textsuperscript{pq} ba for a ∈ A\textsuperscript{p}, b ∈ A\textsuperscript{q}), with a degree-one k-linear derivation d satisfying the graded Leibniz rule and d² = 0. Write A\textsuperscript{0} for the (commutative) degree-0 part. Curvature c ∈ A\textsuperscript{2} with dc = 0 is carried in the definitions because it costs nothing; \textbf{this project only ever uses c = 0} (c ≠ 0 is Z3).

\begin{definition}[A1 — Bounded graded-commutative dga]
\label{A1}
\lean{Cohesive.GCDga}
\leanok
Let k be a commutative ring. A \emph{bounded graded-commutative curved dga} over k is a ℤ-graded k-algebra A = ⨁\textsubscript{q} A\textsuperscript{q} together with

\begin{itemize}
\item a top degree N with A\textsuperscript{q} = 0 for q < 0 and for q > N;
\item graded commutativity: ab = (−1)\textsuperscript{pq} ba for a ∈ A\textsuperscript{p}, b ∈ A\textsuperscript{q};
\item a k-linear differential d of degree +1 with d(ab) = da·b + (−1)\textsuperscript{p} a·db for a ∈ A\textsuperscript{p}, and d ∘ d = 0;
\item a curvature c ∈ A\textsuperscript{2} with dc = 0.
\end{itemize}
Deviation: Block allows non-commutative A with d² = [c, −]; for graded-commutative A this is d² = 0. Boundedness is added (it holds for the Dolbeault algebra).
\par\smallskip{\footnotesize Source: blueprint.md §2.}
\end{definition}
\paragraph{Blueprint notes.}
Definition of the setting above. Encoding: internal ℤ-grading, see conventions §5. Lean: \texttt{GradedAlgebra} ✓; no graded-commutativity class in Mathlib (✗ as a class), so state it as a field. \texttt{class GCDga} in \texttt{Algebra/Dga.lean}, with A\textsuperscript{0} a commutative ring and A an A\textsuperscript{0}-algebra built from graded commutativity. \warn{} ℤ-indexed (zero in negative degrees) rather than ℕ-indexed.

\begin{theorem}[A2 — Connections]
\label{A2}
\lean{Cohesive.Connection, Cohesive.Connection.exists_connection, Cohesive.Connection.diff}
\leanok
Let A be a bounded graded-commutative curved dga over a commutative ring k, and P an A\textsuperscript{0}-module. A \emph{connection} on P is an additive map ∇ : P → A\textsuperscript{1} ⊗\textsubscript{A\textasciicircum{}0} P with ∇(fe) = df ⊗ e + f·∇e for f ∈ A\textsuperscript{0} (it is then automatically k-linear).

\begin{itemize}
\item Every finitely generated projective A\textsuperscript{0}-module has a connection.
\item The difference ∇₁ − ∇₂ of two connections on P is A\textsuperscript{0}-linear.
\end{itemize}
Deviation: additivity is assumed and k-linearity derived; left modules.
\par\smallskip{\footnotesize Source: blueprint.md §2.}
\end{theorem}
\begin{proof}
\leanok
\uses{A1}
See the Lean proof and the cited items.
\paragraph{Blueprint notes.}
A connection on an A\textsuperscript{0}-module P is a k-linear ∇ : P → A\textsuperscript{1} ⊗\textsubscript{A\textasciicircum{}0} P with ∇(fe) = df ⊗ e + f∇e. Extension to A ⊗\textsubscript{A\textasciicircum{}0} P by ∇(a ⊗ e) = da ⊗ e + (−1)\textsuperscript{|a|} a·∇e.

\begin{itemize}
\item \textbf{A2.a} Every finitely generated projective A\textsuperscript{0}-module has a connection. Proof: P ⊕ P' ≅ (A\textsuperscript{0})\textsuperscript{m}; take ∇ = (1 ⊗ π) ∘ d\textsuperscript{⊕m} ∘ ι. [B, proof of Lemma 3.2.8].
\item \textbf{A2.b} The difference of two connections is A\textsuperscript{0}-linear.
\end{itemize}
\end{proof}

\begin{theorem}[A3 — Cohesive and quasi-cohesive modules]
\label{A3}
\lean{Cohesive.QuasiCohesive, Cohesive.CohesiveModule, Cohesive.QuasiCohesive.ofGenerators, Cohesive.QuasiCohesive.X_eq_of_generators, Cohesive.QuasiCohesive.comp_one_leibniz, Cohesive.QuasiCohesive.X0}
\leanok
Let A be a bounded graded-commutative curved dga over k with curvature c. A \emph{quasi-cohesive module} M = (Q, 𝕏) is a family of A\textsuperscript{0}-modules Q\textsuperscript{p} (p ∈ ℤ) together with a k-linear operator 𝕏 of total degree 1 on the total module Tot Q = A ⊗\textsubscript{A\textasciicircum{}0} ⨁\textsubscript{p} Q\textsuperscript{p} (A acting on the left factor) such that 𝕏(a·x) = da·x + (−1)\textsuperscript{|a|} a·𝕏x for homogeneous a, and 𝕏 ∘ 𝕏 = c·(−). It is \emph{cohesive} if Q\textsuperscript{p} = 0 outside a bounded range of p and every Q\textsuperscript{p} is finitely generated projective over A\textsuperscript{0}; P\textsubscript{A} denotes the cohesive modules.

\begin{itemize}
\item 𝕏 is determined by its values on the generators 1 ⊗ e, and conversely an additive map θ on the generators of total degree 1 with θ(fe) = df ⊗ e + f·θ(e) that is flat on generators extends to a quasi-cohesive module with 𝕏(a ⊗ e) = da ⊗ e + (−1)\textsuperscript{|a|} a·θ(e).
\item Writing 𝕏\textsuperscript{j} for the form-degree-j part of 𝕏 on generators: 𝕏\textsuperscript{1}(fe) = df ⊗ e + f·𝕏\textsuperscript{1}(e) (a connection, no sign); 𝕏\textsuperscript{0} is an A\textsuperscript{0}-linear map of degree +1 on ⨁\textsubscript{p} Q\textsuperscript{p} with 𝕏\textsuperscript{0} ∘ 𝕏\textsuperscript{0} = 0.
\end{itemize}
Deviation: left modules; the curvature is c = −c\textsubscript{Block}.
\par\smallskip{\footnotesize Source: blueprint.md §2.}
\end{theorem}
\begin{proof}
\leanok
\uses{A1}
See the Lean proof and the cited items.
\paragraph{Blueprint notes.}
A \emph{quasi-cohesive module} is a ℤ-graded A\textsuperscript{0}-module Q together with a k-linear operator 𝕏 of total degree 1 on A ⊗\textsubscript{A\textasciicircum{}0} Q (left A-module, A acting on the left factor) such that

\begin{itemize}
\item Leibniz: 𝕏(a·m) = da·m + (−1)\textsuperscript{|a|} a·𝕏m, and
\item flatness: 𝕏 ∘ 𝕏 = c·(−) (left multiplication; c = 0 here).
\end{itemize}
It is \emph{cohesive} if moreover Q\textsuperscript{p} ≠ 0 for only finitely many p and each Q\textsuperscript{p} is finitely generated projective. P\textsubscript{A} denotes the cohesive ones. [B, Def 2.3.2] \warn{} left modules, see conventions §3 for the translation to Block's right modules.

\begin{itemize}
\item \textbf{A3.a Components.} Restriction to Q determines 𝕏, with components 𝕏\textsuperscript{j} : Q\textsuperscript{p} → A\textsuperscript{j} ⊗ Q\textsuperscript{p+1−j}; 𝕏\textsuperscript{0} is an A\textsuperscript{0}-linear differential, 𝕏\textsuperscript{1} is a connection on each Q\textsuperscript{p} (no sign in the left convention), 𝕏\textsuperscript{j} (j ≥ 2) are A\textsuperscript{0}-linear. [B, after Def 2.3.1], [A, eq. (2)].
\end{itemize}
A3.b in degree 0 (\texttt{X0\_X0}) is part of this item; the other degrees are A3.b′ (not formalized, not needed).
\end{proof}

\begin{theorem}[A3.b′ — Flatness in all degrees]
\label{A3.bprime}
\uses{A3}
Let M = (Q, 𝕏) be a quasi-cohesive module over a bounded graded-commutative curved dga A with curvature c ∈ A\textsuperscript{2}, and 𝕏\textsuperscript{j} : Q\textsuperscript{p} → A\textsuperscript{j} ⊗\textsubscript{A\textasciicircum{}0} Q\textsuperscript{p+1−j} the components of 𝕏 on generators, extended to Tot Q by the Leibniz rule. Then flatness 𝕏 ∘ 𝕏 = c·(−) is equivalent to the equations Σ\textsubscript{i+j=n} 𝕏\textsuperscript{i} ∘ 𝕏\textsuperscript{j} = δ\textsubscript{n,2}·c·(−) for all n ≥ 0, one in each form degree n (in degree 0: 𝕏\textsuperscript{0} ∘ 𝕏\textsuperscript{0} = 0). Not formalized; not used by any other item.
\par\smallskip{\footnotesize Source: blueprint.md §2.}
\end{theorem}
\begin{proof}
See the Lean proof and the cited items.
\paragraph{Blueprint notes.}
\textbf{A3.b} Flatness unfolds degree by degree into the equations of [A, eq. (2)].
\end{proof}

\begin{definition}[A4 — Hom complexes]
\label{A4}
\lean{Cohesive.Hom, Cohesive.δ, Cohesive.homEquivGen}
\leanok
\uses{A3}
Let A be a bounded graded-commutative curved dga over k and M₁ = (Q₁, 𝕏₁), M₂ = (Q₂, 𝕏₂) quasi-cohesive modules. Hom\textsuperscript{n}(M₁, M₂) is the k-module of k-linear maps φ : Tot Q₁ → Tot Q₂ raising total degree by n with φ(a·x) = (−1)\textsuperscript{n|a|} a·φ(x) for homogeneous a, with differential δφ = 𝕏₂ ∘ φ − (−1)\textsuperscript{n} φ ∘ 𝕏₁.

\begin{itemize}
\item δ maps Hom\textsuperscript{n} to Hom\textsuperscript{n+1} and δ ∘ δ = 0 (also with curvature, since |c| = 2).
\item Composition maps Hom\textsuperscript{m} × Hom\textsuperscript{n} to Hom\textsuperscript{m+n}, and δ(ψ ∘ φ) = δψ ∘ φ + (−1)\textsuperscript{m} ψ ∘ δφ for ψ of degree m; the identity lies in Hom\textsuperscript{0} and is closed.
\item δ agrees with Mathlib's \texttt{CochainComplex.HomComplex.δ}: δφ = 𝕏₂ ∘ φ + (−1)\textsuperscript{n+1} φ ∘ 𝕏₁.
\item Restriction to generators is a k-linear isomorphism from Hom\textsuperscript{n}(M₁, M₂) to the A\textsuperscript{0}-linear maps ψ : ⨁\textsubscript{p} Q₁\textsuperscript{p} → Tot Q₂ with ψ(Q₁\textsuperscript{p}) in total degree p + n; the inverse is ψ̂(a ⊗ e) = (−1)\textsuperscript{n|a|} a·ψ(e).
\end{itemize}
Deviation: M₁ may be any quasi-cohesive module (the sources take M₁ cohesive).
\par\smallskip{\footnotesize Source: blueprint.md §2.}
\end{definition}
\paragraph{Blueprint notes.}
For M₁ cohesive and M₂ quasi-cohesive: Hom\textsuperscript{n}(M₁, M₂) = k-linear maps φ of degree n with φ(a·m) = (−1)\textsuperscript{n|a|} a·φ(m); differential δφ = 𝕏₂∘φ − (−1)\textsuperscript{n} φ∘𝕏₁.

\begin{itemize}
\item \textbf{A4.a} δ∘δ = 0, including with curvature (|c| even). [B, Lemma 2.3.4]
\item \textbf{A4.b} Composition is associative, unital and satisfies the Leibniz rule for δ.
\item \textbf{A4.c} Restriction to generators gives Hom\textsuperscript{n}(M₁, M₂) ≅ ⊕\textsubscript{p} Hom\textsubscript{A\textasciicircum{}0}(Q₁\textsuperscript{p}, (A ⊗ Q₂)\textsuperscript{p+n}).
\item \textbf{A4.d} The sign convention agrees with Mathlib's \texttt{CochainComplex.HomComplex.δ} ✓ (δz = z ≫ d + (−1)\textsuperscript{n+1} d ≫ z), conventions §2.
\end{itemize}

\begin{definition}[A5 — Homotopy category Ho]
\label{A5}
\lean{Cohesive.HoP}
\leanok
\uses{A4}
Let A be a bounded graded-commutative curved dga over k. The homotopy category Ho(P\textsubscript{A}) has the cohesive modules as objects and, as morphisms M → N, the degree-0 cohomology H\textsuperscript{0} = Z\textsuperscript{0}/B\textsuperscript{0} of the Hom complex, where Z\textsuperscript{0} = \{φ ∈ Hom\textsuperscript{0}(M, N) : δφ = 0\} and B\textsuperscript{0} = δ(Hom\textsuperscript{−1}(M, N)); composition is induced by composition of maps. Ho(P\textsubscript{A}) is a preadditive, k-linear category.
\par\smallskip{\footnotesize Source: blueprint.md §2.}
\end{definition}
\paragraph{Blueprint notes.}
Objects: cohesive modules. Morphisms: H\textsuperscript{0} of the Hom complex. A category (from A4.b), k-linear. Homotopy equivalences = isomorphisms in Ho(P\textsubscript{A}). [B, §2.4]

\begin{theorem}[A6 — Form-degree filtration]
\label{A6}
\lean{Cohesive.filTot, Cohesive.filHom, Cohesive.memFilHom_iff, Cohesive.δ_mem_filHom}
\leanok
Let A be a bounded graded-commutative curved dga over k with top degree N, and M₁, M₂ quasi-cohesive modules. F\textsuperscript{b} Tot Q is the k-span of the a ⊗ e with a ∈ A\textsuperscript{q}, q ≥ b, and F\textsuperscript{b} Hom\textsuperscript{n}(M₁, M₂) consists of the φ ∈ Hom\textsuperscript{n} with φ(F\textsuperscript{c} Tot Q₁) ⊆ F\textsuperscript{c+b} Tot Q₂ for all c.

\begin{itemize}
\item φ ∈ Hom\textsuperscript{n} lies in F\textsuperscript{b} iff φ(1 ⊗ e) ∈ F\textsuperscript{b} Tot Q₂ for every generator e, i.e. the components of φ into A\textsuperscript{b′} ⊗ Q₂ vanish for b′ < b.
\item δ maps F\textsuperscript{b} Hom\textsuperscript{n} into F\textsuperscript{b} Hom\textsuperscript{n+1}.
\item The filtration is finite: F\textsuperscript{0} Hom\textsuperscript{n} = Hom\textsuperscript{n} and F\textsuperscript{b} Hom\textsuperscript{n} = 0 for b > N.
\end{itemize}
\par\smallskip{\footnotesize Source: blueprint.md §2.}
\end{theorem}
\begin{proof}
\leanok
\uses{A4}
See the Lean proof and the cited items.
\paragraph{Blueprint notes.}
F\textsuperscript{b} Hom\textsuperscript{n} = maps whose components into A\textsuperscript{b'} ⊗ − vanish for b' < b. δ preserves it, the filtration is finite (0 ≤ b ≤ N), and gr\textsuperscript{b} ≅ Hom\textsubscript{A\textasciicircum{}0}(Q₁, A\textsuperscript{b} ⊗ Q₂) with differential induced by 𝕏₁\textsuperscript{0}, 𝕏₂\textsuperscript{0}. [B, Prop 2.5.1], [A, Def 3.3.2]
\end{proof}

\begin{theorem}[A7 — Finite filtration comparison]
\label{A7}
\lean{Cohesive.quasiIso_of_finite_filtration}
\leanok
Let C be an abelian category and N ∈ ℕ. Suppose given, for b = 0, …, N, maps f\textsubscript{b} : F\textsuperscript{b} K → F\textsuperscript{b} L of ℤ-indexed cochain complexes in C, and for b < N maps of short exact sequences 0 → F\textsuperscript{b+1} → F\textsuperscript{b} → gr\textsuperscript{b} → 0 (of the K-side into the L-side) whose first two components are f\textsubscript{b+1} and f\textsubscript{b}. If every induced map gr\textsuperscript{b} K → gr\textsuperscript{b} L (b < N) is a quasi-isomorphism and f\textsubscript{N} is a quasi-isomorphism (e.g. F\textsuperscript{N} = 0), then f\textsubscript{0} is a quasi-isomorphism.

Deviation: replaces the spectral sequences of the sources by the homology long exact sequence.
\par\smallskip{\footnotesize Source: blueprint.md §2.}
\end{theorem}
\begin{proof}
\leanok
See the Lean proof and the cited items.
\paragraph{Blueprint notes.}
A map of finitely filtered complexes inducing quasi-isomorphisms on all gr\textsuperscript{b} is a quasi-isomorphism. Proof: induction on the filtration length with the long exact sequence. Replaces the spectral sequences in [B] and [A]. Lean: homology long exact sequence ✓ (\texttt{HomologicalComplex.HomologySequence}).
\end{proof}

\begin{theorem}[A8 — Homotopy criterion]
\label{A8}
\uses{A7, F15, Z2.g}
Let A be a bounded graded-commutative dga over k and φ : M₁ → M₂ a closed degree-0 map of cohesive modules. Then φ is an isomorphism in Ho(P\textsubscript{A}) iff its (0,0)-component φ\textsuperscript{0} : (Q₁, 𝕏₁\textsuperscript{0}) → (Q₂, 𝕏₂\textsuperscript{0}) is a homotopy equivalence of complexes of A\textsuperscript{0}-modules (equivalently, since these are bounded complexes of projectives, a quasi-isomorphism). Not formalized; not used by any other item.
\par\smallskip{\footnotesize Source: blueprint.md §2.}
\end{theorem}
\begin{proof}
See the Lean proof and the cited items.
\paragraph{Blueprint notes.}
Not on the critical path: no step of C or F uses it. Schedule it after F15, or with Z2. A closed degree-0 φ : M₁ → M₂ between cohesive modules is an isomorphism in Ho(P\textsubscript{A}) iff φ\textsuperscript{0} : (Q₁, 𝕏₁\textsuperscript{0}) → (Q₂, 𝕏₂\textsuperscript{0}) is a homotopy equivalence of complexes of A\textsuperscript{0}-modules. [A, Lemma 3.3.5], [B, Prop 2.5.2]

\begin{itemize}
\item Proof (⇐): for every cohesive M, φ\textsubscript{\emph{} : Hom(M, M₁) → Hom(M, M₂) induces on gr\textsuperscript{b} the map Hom\textsubscript{A\textasciicircum{}0}(Q, A\textsuperscript{b} ⊗ φ\textsuperscript{0}), a homotopy equivalence (any additive functor preserves them); by A7 φ\textsubscript{}} is a quasi-isomorphism; Yoneda as in [A, Lemma 3.3.5].
\item Proof (⇒): take the (0,0)-components of the homotopies.
\item \textbf{A8.a} For bounded complexes of projective A\textsuperscript{0}-modules, homotopy equivalence ⇔ quasi-isomorphism ([B]'s form). Lean: \texttt{CochainComplex.isKProjective\_of\_projective} ✓.
\end{itemize}
\end{proof}

\begin{theorem}[A9 — Base change along dga maps]
\label{A9}
\lean{Cohesive.DgaHom, Cohesive.CohesiveModule.baseChange, Cohesive.bcHoFunctor, Cohesive.bcCompNatIso, Cohesive.bcHomL, Cohesive.hom_eq_of_bcF, Cohesive.DgaHom.comp, Cohesive.bcCompTot}
\leanok
Let A, B be bounded graded-commutative curved dgas over k. A \emph{dga map} f : A → B is a k-algebra map preserving degrees, the differential and the curvature (f(c\textsubscript{A}) = c\textsubscript{B}). For a quasi-cohesive A-module M = (Q, 𝕏), the base change f\textsuperscript{*M} is the quasi-cohesive B-module with generators B\textsuperscript{0} ⊗\textsubscript{A\textasciicircum{}0} Q\textsuperscript{p} and 𝕏\textsubscript{B}(b ⊗ m) = db ⊗ m + (−1)\textsuperscript{|b|} b·𝕏m.

\begin{itemize}
\item If M is cohesive, so is f\textsuperscript{*M}.
\item φ ∈ Hom\textsuperscript{n}(M₁, M₂) gives f\textsuperscript{\emph{φ} ∈ Hom\textsuperscript{n}(f\textsuperscript{*M₁}, f\textsuperscript{*M₂}), (f\textsuperscript{*φ})(b ⊗ m) = (−1)\textsuperscript{n|b|} b·φ(m), with f\textsuperscript{}}(δφ) = δ(f\textsuperscript{*φ}); this gives a functor Ho(P\textsubscript{A}) ⥤ Ho(P\textsubscript{B}).
\item Maps of degree n out of f\textsuperscript{*M₁} are determined by their values on the image of Tot Q₁ → Tot(f\textsuperscript{*Q₁}), a ⊗ e ↦ f(a) ⊗ (1 ⊗ e).
\item For a second dga map g : B → C, γ ⊗ q ↦ γ ⊗ (1 ⊗ q) is a closed degree-0 isomorphism (g ∘ f)\textsuperscript{\emph{M} ≅ g\textsuperscript{}}(f\textsuperscript{\emph{M}), natural on Hom complexes in every degree; it gives a natural isomorphism of functors (g ∘ f)\textsuperscript{}} ≅ g\textsuperscript{\emph{} ∘ f\textsuperscript{}} on Ho(P\textsubscript{A}).
\end{itemize}
Deviation: left modules; Block's ω is 0.
\par\smallskip{\footnotesize Source: blueprint.md §2.}
\end{theorem}
\begin{proof}
\leanok
\uses{A5}
See the Lean proof and the cited items.
\paragraph{Blueprint notes.}
For a dga map f : A → B, M = (Q, 𝕏) ↦ f\textsuperscript{*M} with generators B\textsuperscript{0} ⊗\textsubscript{A\textasciicircum{}0} Q\textsuperscript{p} (Block's E\textsuperscript{•} ⊗\textsubscript{𝒜\_1} 𝒜\textsubscript{2}, [B, §2.7]) and 𝕏\textsubscript{B}(b ⊗ m) = db ⊗ m + (−1)\textsuperscript{|b|} b·𝕏m; its total module B ⊗\textsubscript{B\textasciicircum{}0} ⨁\textsubscript{p} (B\textsuperscript{0} ⊗\textsubscript{A\textasciicircum{}0} Q\textsuperscript{p}) is canonically B ⊗\textsubscript{A} M. Cohesive goes to cohesive (B\textsuperscript{0} ⊗\textsubscript{A\textasciicircum{}0} Q\textsuperscript{p} is finitely generated projective); a dg functor; preserves homotopy equivalences. Used for restriction of cohesive modules to open subsets. [B, §2.7 with ω = 0]

\begin{itemize}
\item \textbf{A9.a} Base change along a composite B → C is naturally isomorphic to the composite of base changes, compatibly with Hom complexes (used when restricting from X to U′ to V ⊆ U′).
\end{itemize}
\end{proof}

\begin{theorem}[A10 — Cohesive modules from complexes]
\label{A10}
\lean{Cohesive.QuasiCohesive.fromComplex, Cohesive.IsBoundedFGProj, Cohesive.fcResP_δ, Cohesive.homRδ, Cohesive.homRδ_mem, Cohesive.fc_hom_ext}
\leanok
Let A be a bounded graded-commutative dga over k with curvature c = 0, ρ : R → A\textsuperscript{0} a map of commutative rings with d ∘ ρ = 0, and (W, d\textsubscript{W}) a cochain complex of R-modules. Then A ⊗\textsubscript{R} W, with generators A\textsuperscript{0} ⊗\textsubscript{R} W\textsuperscript{p} and 𝕏(a ⊗ w) = da ⊗ w + (−1)\textsuperscript{|a|} a ⊗ d\textsubscript{W} w, is a quasi-cohesive module; it is cohesive if W is bounded with finitely generated projective terms. For every quasi-cohesive N, restriction along w ↦ 1 ⊗ w is an additive bijection from Hom\textsuperscript{n}(A ⊗\textsubscript{R} W, N) to the families of ρ-semilinear maps f\textsubscript{p} : W\textsuperscript{p} → Tot N landing in total degree p + n; it carries δ to (δf)\textsubscript{p} = 𝕏\textsubscript{N} ∘ f\textsubscript{p} − (−1)\textsuperscript{n} f\textsubscript{p+1} ∘ d\textsubscript{W}.

Deviation: requires c = 0 (otherwise 𝕏 ∘ 𝕏 = 0 ≠ c·(−)).
\par\smallskip{\footnotesize Source: blueprint.md §2.}
\end{theorem}
\begin{proof}
\leanok
\uses{A9}
See the Lean proof and the cited items.
\paragraph{Blueprint notes.}
Let ρ : R → A\textsuperscript{0} be a map of commutative rings with d∘ρ = 0, and (W, d\textsubscript{W}) a bounded complex of finitely generated projective R-modules. Then A ⊗\textsubscript{R} W with 𝕏 = d ⊗ 1 + 1 ⊗ d\textsubscript{W} is cohesive. (v1.10) With the Koszul sign of conventions §2: 𝕏(a ⊗ w) = da ⊗ w + (−1)\textsuperscript{|a|} a ⊗ d\textsubscript{W} w; on generators 1 ⊗ w no sign appears, and this is what \texttt{fromComplex} implements.

\begin{itemize}
\item \textbf{A10.a Adjunction.} Hom\textsuperscript{•\_}\{P\textsubscript{A}\}(A ⊗\textsubscript{R} W, M) ≅ Hom\textsuperscript{•\_R}(W, M) as complexes, via precomposition with w ↦ 1 ⊗ w. Used in C4.
\end{itemize}
\end{proof}

\begin{theorem}[A11 — Lifting theorem]
\label{A11}
\lean{Cohesive.lifting, Cohesive.CohesiveModule.ofGenData, Cohesive.lift_core, Cohesive.liftStage_one, Cohesive.liftStage_succ, Cohesive.ConeLift}
\leanok
\uses{A2}
Let A be a bounded graded-commutative curved dga over k (curvature c arbitrary) such that every A\textsuperscript{b} is flat over A\textsuperscript{0}. Let N = (Q, 𝕏) be a quasi-cohesive module, (P, ∂) a bounded cochain complex of finitely generated projective A\textsuperscript{0}-modules, and e\textsuperscript{0} : (P, ∂) → (Q, 𝕏\textsuperscript{0}) a quasi-isomorphism. Then there are a cohesive module (A ⊗\textsubscript{A\textasciicircum{}0} P, 𝔼) with generators P\textsuperscript{p} and 𝔼\textsuperscript{0} = ∂, and a closed degree-0 map e : (A ⊗ P, 𝔼) → N whose (0,0)-component is e\textsuperscript{0} (ε(e(1 ⊗ x)) = e\textsuperscript{0}(x), ε the form-degree-0 part). The proof constructs 𝔼 and e form degree by form degree; at each step it only needs that the cone of e\textsuperscript{0} ∘ − : gr\textsuperscript{b} Hom(P, P) → gr\textsuperscript{b} Hom(P, Q) has no cohomology in degree 1.
\par\smallskip{\footnotesize Source: blueprint.md §2.}
\end{theorem}
\begin{proof}
\leanok
\uses{A6}
See the Lean proof and the cited items.
\paragraph{Blueprint notes.}
Assume each A\textsuperscript{b} is flat over A\textsuperscript{0}. Let N = (Q, 𝕏) be quasi-cohesive, (P, ∂) a bounded complex of finitely generated projective A\textsuperscript{0}-modules, and e\textsuperscript{0} : (P, ∂) → (Q, 𝕏\textsuperscript{0}) a quasi-isomorphism. Then there are a cohesive structure 𝔼 on A ⊗ P with 𝔼\textsuperscript{0} = ∂ and a closed degree-0 e : (A ⊗ P, 𝔼) → N with e\textsuperscript{0} as (0,0)-component. [B, Thm 3.2.7(2) and Lemma 3.2.8], [A, Prop 3.3.4, case 2]

\begin{itemize}
\item Proof: choose connections on each P\textsuperscript{p} (A2.a). Construct 𝔼\textsuperscript{b} and e\textsuperscript{b} by induction on b. At each step the defect pair (P\textsubscript{b}, Q\textsubscript{b}) of [A, (10)–(12)] is a cycle in the cone of the map of A11.a, which is acyclic; a preimage gives (𝔼\textsuperscript{b+1}, e\textsuperscript{b+1}). Stop at b = N.
\item \warn{} Re-derive every sign of [A, (10)–(12)] in the left convention before coding (Antweiler composes with a Koszul sign (−1)\textsuperscript{a'b}, conventions §4).
\item \textbf{A11.a Internal hypothesis.} The proof only uses that, for every b, the map Hom\textsubscript{A\textasciicircum{}0}(P, A\textsuperscript{b} ⊗ P) → Hom\textsubscript{A\textasciicircum{}0}(P, A\textsuperscript{b} ⊗ Q) induced by e\textsuperscript{0} is a quasi-isomorphism ([A]'s form). This follows from the assumptions: A\textsuperscript{b} ⊗ − preserves quasi-isomorphisms, and so does Hom out of a bounded complex of projectives.
\end{itemize}
\end{proof}

\begin{theorem}[A12 — Total quasi-isomorphism]
\label{A12}
\lean{Cohesive.totMap_quasiIso, Cohesive.lifting_quasiIso, Cohesive.QuasiCohesive.totComplex, Cohesive.totMap, Cohesive.Z0.comp0}
\leanok
\uses{A11, A7}
Let A be a bounded graded-commutative dga over k with curvature c = 0 such that every A\textsuperscript{b} is flat over A\textsuperscript{0}. Let φ : M₁ → M₂ be a closed degree-0 map of quasi-cohesive modules whose (0,0)-component φ\textsuperscript{0} : (Q₁, 𝕏₁\textsuperscript{0}) → (Q₂, 𝕏₂\textsuperscript{0}) is a quasi-isomorphism of complexes of A\textsuperscript{0}-modules. Then φ is a quasi-isomorphism of total complexes (Tot Q₁, 𝕏₁) → (Tot Q₂, 𝕏₂) (complexes of k-modules graded by total degree). In particular, in the lifting theorem (A11) with c = 0, the map e is a quasi-isomorphism of total complexes.

Deviation: stated for any such φ (the sources state it for the map e of A11).
\par\smallskip{\footnotesize Source: blueprint.md §2.}
\end{theorem}
\begin{proof}
\leanok
See the Lean proof and the cited items.
\paragraph{Blueprint notes.}
In A11, e is a quasi-isomorphism of total complexes. Proof: A7 with the form-degree filtration. [A, proof of Prop 7.2.18]
\end{proof}

\begin{definition}[A13 — Bridge to Block's definition]
\label{A13}
\lean{Cohesive.TotR.rAct, Cohesive.BlockCohesive, Cohesive.BlockHom, Cohesive.blockδ, Cohesive.blockSwap, Cohesive.BlockQuasiCohesive.toLeft}
\leanok
\uses{A4}
Let A be a bounded graded-commutative curved dga over k with curvature c, and c\textsubscript{B} ∈ A with c = −c\textsubscript{B}. A quasi-cohesive module in Block's convention is a family of A\textsuperscript{0}-modules Q\textsuperscript{p} with a k-linear operator 𝔼 of total degree 1 on the right module (⨁\textsubscript{p} Q\textsuperscript{p}) ⊗\textsubscript{A\textasciicircum{}0} A satisfying 𝔼(eω) = 𝔼(e ⊗ 1)·ω + (−1)\textsuperscript{|e|} e ⊗ dω for e ∈ Q\textsuperscript{p}, and 𝔼 ∘ 𝔼 + (−)·c\textsubscript{B} = 0; its morphisms of degree n are the right A-linear maps of degree n, with dφ = 𝔼₂ ∘ φ − (−1)\textsuperscript{n} φ ∘ 𝔼₁. Transport along σ : e ⊗ a ↦ (−1)\textsuperscript{|e||a|} a ⊗ e gives

\begin{itemize}
\item mutually inverse bijections between Block's (cohesive, quasi-cohesive) modules and those of A3;
\item for each n a k-linear isomorphism from Block's Hom\textsuperscript{n} to Hom\textsuperscript{n} of A4, φ ↦ σ ∘ φ ∘ σ\textsuperscript{−1}, compatible with the differentials, composition and identities.
\end{itemize}
\par\smallskip{\footnotesize Source: blueprint.md §2.}
\end{definition}
\paragraph{Blueprint notes.}
The left-module definition is dg-equivalent to [B, Def 2.3.2] (right modules) for graded-commutative A, with curvature c ↔ −c. Not on the critical path; proves the conventions are faithful.
